\documentclass[a4paper,12pt]{article}

\usepackage[a4paper,margin=2.5cm]{geometry}
\usepackage{iftex}
\ifPDFTeX
  \usepackage[T2A]{fontenc}
  \usepackage[utf8]{inputenc}
  \usepackage[russian]{babel}
\else
  \usepackage[russian]{babel}
  \usepackage{fontspec}
  \setmainfont{Times New Roman}
\fi

\usepackage{amsmath}
\usepackage{amssymb}
\usepackage{enumitem}

\begin{document}

\begin{center}
{\LARGE Коллоквиум по курсу}\\[0.2em]
{\LARGE «Уравнения математической физики»}\\[1em]
{\large лекции 1--7: введение, модели, классификация, корректность}
\end{center}

\vspace{1.2em}
\section*{Часть 1. Теоретические вопросы}

\begin{enumerate}
\item \textbf{Математическое моделирование}
  \begin{enumerate}[label=(\alph*)]
  \item Перечислите основные этапы построения математической модели физического процесса.

  **Для себя**
  Мат. моделирование - метод изучения реальных объектов , процессов или явлений путем построения их математических "образов"/ моделей и последуещего анализа этих моделей с помощью мат методов.


Выделим основные этапы, опираясь на классическую схему:

\begin{enumerate}
    \item \textbf{Выбор величин, описывающих процесс.} Обычно это функции координат пространства 
    $x = (x_1, \dots, x_n)$ и времени $t$: например, плотность $\rho(x,t)$, температура $T(x,t)$, 
    скорость $\mathbf{u}(x,t)$ и т.д.

    \item \textbf{Формулировка законов, которым подчиняется идеализированный процесс.} 
    На основе фундаментальных законов физики (сохранения массы, энергии, импульса и т.п.) 
    выводятся математические соотношения — чаще всего дифференциальные, интегральные 
    или функциональные уравнения.

    \textit{Сплошная среда} представляет собой материальный континуум, то есть непрерывное 
    множество материальных точек с непрерывным распределением по нему кинематических, 
    динамических, термодинамических и иных физико-химических характеристик рассматриваемой среды.

    Явления, рассматриваемые в механике сплошных сред, носят макроскопический характер. 
    Это позволяет абстрагироваться от молекулярного строения вещества и рассматривать 
    физические тела как сплошные среды. Понятие «сплошная среда» представляет собой 
    модель реальных сред.

    \item \textbf{Задание дополнительных условий.} Поскольку дифференциальные уравнения 
    обычно имеют бесконечно много решений, для выделения конкретного процесса необходимо 
    задать начальные и граничные условия.

    \item \textbf{Исследование корректности задачи.} Проверяют, существует ли решение, 
    единственно ли оно и устойчиво ли (непрерывно зависит от исходных данных).

    \item \textbf{Нахождение решения.} Используются аналитические или численные методы, 
    ориентированные на ЭВМ.

    \item \textbf{Анализ свойств решения и физическая интерпретация.} 
    На основе полученного решения делаются выводы о протекании исследуемого процесса.
\end{enumerate}
  \item Какую роль играет лемма 1.1 (об обращении интеграла в нуль) при выводе дифференциальных уравнений?
  
  \textbf{Лемма 1.1.} \textit{Пусть $D$ — произвольная область с границей $S$, 
и пусть непрерывная в $D$ функция $\psi$ обладает тем свойством, что для любой 
ограниченной (кубируемой) области $\Omega \subset D$}
\begin{equation}
\int\limits_{\Omega} \psi(\mathbf{x})\,dx = 0.
\end{equation}

\textit{Тогда $\psi(\mathbf{x}) \equiv 0$ в $D$.}

\textbf{Доказательство.} Предположим противное, что $\psi(\mathbf{x}) \ne 0$. 
Тогда найдется такая точка $\mathbf{x}_0 \in D$, что $\psi(\mathbf{x}_0) = \varepsilon \ne 0$. 
Пусть, например, $\varepsilon > 0$. В силу непрерывности $\psi$ существует такая окрестность 
$U_\delta = U_\delta(\mathbf{x}_0)$ с центром в точке $\mathbf{x}_0$ радиуса $\delta$, 
что $\psi(\mathbf{x}) \ge \varepsilon/2$ для всех $\mathbf{x} \in U_\delta$. 
Полагая $\Omega = U_\delta(\mathbf{x}_0)$, имеем в силу (1.1):
\[
0 = \int\limits_{U_\delta} \psi(\mathbf{x})\,dx 
\ge \frac{\varepsilon}{2} \int\limits_{U_\delta} dx 
= \frac{2\varepsilon}{3}\pi \delta^3 > 0.
\]

Полученное противоречие доказывает лемму. $\blacksquare$
  \end{enumerate}

\item \textbf{Дифференциальные операторы}
  \begin{enumerate}[label=(\alph*)]
  \item Дайте определение градиента, дивергенции, ротора. Запишите их в декартовых координатах.
  \section{Градиент скалярного поля}

Пусть $u = u(x,y,z)$ — скалярная функция (например, температура, потенциал). 
Градиентом $\mathrm{grad}\,u$ называется вектор, указывающий направление 
наискорейшего возрастания функции и равный по величине скорости этого возрастания.

В декартовой системе координат:
\[
\mathrm{grad}\,u = \frac{\partial u}{\partial x}\,\mathbf{i}
+ \frac{\partial u}{\partial y}\,\mathbf{j}
+ \frac{\partial u}{\partial z}\,\mathbf{k}.
\]

\textbf{Пример 1.} Если $u$ — электростатический потенциал, то вектор напряженности 
электрического поля связан с градиентом соотношением 
$\mathbf{E} = -\mathrm{grad}\,u$. 
В гравитации $\mathbf{f} = \mathrm{grad}\,u$ определяет силу притяжения на единичную массу.

\section{Дивергенция векторного поля}

Для векторного поля $\mathbf{v} = (v_x, v_y, v_z)$ дивергенция определяется как
\[
\mathrm{div}\,\mathbf{v} = \frac{\partial v_x}{\partial x}
+ \frac{\partial v_y}{\partial y}
+ \frac{\partial v_z}{\partial z}.
\]

Физический смысл: дивергенция характеризует мощность источников поля. 
Если $\mathrm{div}\,\mathbf{v} > 0$ в точке, то из нее исходит поток; 
если $\mathrm{div}\,\mathbf{v} < 0$ — поток втекает. 
В гидродинамике уравнение неразрывности 
\[
\mathrm{div}(\rho \mathbf{u}) = 0
\]
выражает закон сохранения массы для несжимаемой жидкости в стационарном случае.

Физический смысл дивергенции векторного поля — это мера плотности источников 
или стоков в данной точке, показывающая, насколько поток поля, выходящий из 
малой окрестности, превосходит входящий.

Рассмотрим примеры векторных полей и найдем их дивергенцию.

\section{Ротор векторного поля}

Ротор (вихрь) векторного поля определяется формулой:
\[
\mathrm{rot}\,\mathbf{v} =
\begin{vmatrix}
\mathbf{i} & \mathbf{j} & \mathbf{k} \\
\frac{\partial}{\partial x} & \frac{\partial}{\partial y} & \frac{\partial}{\partial z} \\
v_x & v_y & v_z
\end{vmatrix}
=
\left(\frac{\partial v_z}{\partial y} - \frac{\partial v_y}{\partial z}\right)\mathbf{i}
+
\left(\frac{\partial v_x}{\partial z} - \frac{\partial v_z}{\partial x}\right)\mathbf{j}
+
\left(\frac{\partial v_y}{\partial x} - \frac{\partial v_x}{\partial y}\right)\mathbf{k}.
\]

Ротор измеряет завихренность поля. Например, в гидродинамике 
$\mathrm{rot}\,\mathbf{u}$ связан с угловой скоростью вращения частиц жидкости. 
Электромагнитные уравнения Максвелла содержат операторы ротора 
для электрического и магнитного полей.

Рассмотрим пример в случае двух переменных.

\begin{enumerate}
    \item \textbf{Положительный ротор (циркуляция против часовой стрелки).} 
    $\mathbf{v} = \{-y;\, x\}$,
    \[
    \mathrm{rot}\,\mathbf{v} = 1 - (-1) = 2 > 0.
    \]

    \item \textbf{Отрицательный ротор (циркуляция по часовой стрелке).} 
    $\mathbf{v} = \{y;\, -x\}$,
    \[
    \mathrm{rot}\,\mathbf{v} = -1 - 1 = -2 < 0.
    \]

    \item \textbf{Нулевой ротор (безвихревое поле, потенциальное).} 
    $\mathbf{v} = \{x;\, y\}$,
    \[
    \mathrm{rot}\,\mathbf{v} = 0.
    \]
\end{enumerate}
  \item Что такое оператор Лапласа? Приведите его выражение в декартовых, цилиндрических и сферических координатах.\section{Оператор Лапласа (лапласиан)}

Оператор Лапласа — это дифференциальный оператор второго порядка, 
определяемый как дивергенция градиента:
\[
\Delta u = \mathrm{div}(\mathrm{grad}\,u).
\]

\subsection{Выражение в различных системах координат}

\begin{itemize}
    \item \textbf{Декартовы координаты $(x,y,z)$:}
    \[
    \Delta u = \frac{\partial^2 u}{\partial x^2}
    + \frac{\partial^2 u}{\partial y^2}
    + \frac{\partial^2 u}{\partial z^2}.
    \]

    \item \textbf{Цилиндрические координаты $(\rho,\varphi,z)$:}
    \[
    \Delta u =
    \frac{1}{\rho}\frac{\partial}{\partial \rho}
    \left(\rho \frac{\partial u}{\partial \rho}\right)
    + \frac{1}{\rho^2}\frac{\partial^2 u}{\partial \varphi^2}
    + \frac{\partial^2 u}{\partial z^2}.
    \]

    \item \textbf{Сферические координаты $(r,\theta,\varphi)$:}
    \[
    \Delta u =
    \frac{1}{r^2}\frac{\partial}{\partial r}
    \left(r^2 \frac{\partial u}{\partial r}\right)
    + \frac{1}{r^2 \sin\theta}\frac{\partial}{\partial \theta}
    \left(\sin\theta \frac{\partial u}{\partial \theta}\right)
    + \frac{1}{r^2 \sin^2\theta}\frac{\partial^2 u}{\partial \varphi^2}.
    \]
\end{itemize}

\subsection{Физический смысл лапласиана}

В каждой точке лапласиан можно интерпретировать как меру отклонения 
значения функции в этой точке от её среднего значения по малой окрестности. 
Если $\Delta u > 0$, то значение функции меньше среднего, 
если $\Delta u < 0$ — больше среднего.

Именно поэтому уравнение Лапласа
\[
\Delta u = 0
\]
описывает гармонические функции, для которых в каждой точке значение 
равно среднему по любой сфере (в трёхмерном случае) или окружности (на плоскости).

Рассмотрим пример задачи плоского конденсатора с двумя пластинами 
с разноимёнными зарядами.
  \item Докажите тождество \(\operatorname{rot}\,\operatorname{grad} u = 0\).
  \end{enumerate}

\item \textbf{Модели гравитационного и электростатического полей}
  \begin{enumerate}[label=(\alph*)]
  \item Запишите закон всемирного тяготения и закон Кулона. Как вводится потенциал?\begin{itemize}
    \item \textbf{Закон всемирного тяготения и закон Кулона}
    
    Согласно \textit{Замечанию 3.1} на ваших изображениях, сила гравитационного притяжения между двумя точечными массами $m_1$ и $m_2$ определяется как:
    \begin{equation}
    f_{gr} = \gamma \frac{m_1 m_2}{r^2} \tag{3.13}
    \end{equation}
    где $\gamma$ — гравитационная постоянная.
    
    Закон Кулона для двух точечных зарядов $q_1$ и $q_2$ записывается аналогично:
    \begin{equation}
    f_{el} = k \frac{q_1 q_2}{r^2}
    \end{equation}
    где $k$ — электростатическая постоянная.

    \item \textbf{Как вводится потенциал?}
    \begin{itemize}
    \item \textbf{Введение гравитационного потенциала}
    
    Согласно гипотезе дальнодействия Лапласа, введение потенциала опирается на следующие положения:

    \begin{enumerate}
        \item \textbf{Определение интенсивности субстанции:} 
        Наличие любого притягивающего тела массой $m$ порождает во всем пространстве некоторую «субстанцию». Интенсивность этой субстанции $u$ в произвольной точке $\mathbf{x} = (x, y, z)$ вводится формулой:
        \begin{equation}
        u(\mathbf{x}) = \gamma \frac{m}{|\mathbf{x} - \mathbf{x}_0|} \tag{3.1}
        \end{equation}
        где $\mathbf{x}_0$ — точка расположения тела, а $|\mathbf{x} - \mathbf{x}_0|$ — расстояние между точками, вычисляемое как:
        \begin{equation}
        |\mathbf{x} - \mathbf{x}_0| = \sqrt{(x - x_0)^2 + (y - y_0)^2 + (z - z_0)^2} \tag{3.2}
        \end{equation}

        \item \textbf{Связь с вектором силы (Градиент):} 
        Потенциал вводится таким образом, чтобы через него можно было выразить вектор силы притяжения $\mathbf{f}$, действующей на тело единичной массы. Связь устанавливается через операцию градиента:
        \begin{equation}
        \mathbf{f} = \text{grad} u \tag{3.3}
        \end{equation}
        В декартовых координатах это эквивалентно трем скалярным равенствам для компонент силы:
        \begin{equation}
        f_x = \frac{\partial u}{\partial x}, \quad f_y = \frac{\partial u}{\partial y}, \quad f_z = \frac{\partial u}{\partial z} \tag{3.4}
        \end{equation}

        \item \textbf{Принцип суперпозиции:} 
        Для системы из $N$ материальных точек с массами $m_j$, расположенных в точках $\mathbf{x}_j$, суммарный гравитационный потенциал вводится как сумма потенциалов, создаваемых каждым из этих тел в отдельности:
        \begin{equation}
        u(\mathbf{x}) = \gamma \sum_{j=1}^{N} \frac{m_j}{|\mathbf{x} - \mathbf{x}_j|} \tag{3.5}
        \end{equation}
    \end{enumerate}

    \item \textbf{Физический смысл:} 
    Функцию $u$ называют \textit{гравитационным потенциалом}, а вектор $\text{grad} u$ указывает направление быстрейшего возрастания этого потенциала в пространстве.
\end{itemize}
    На странице 30 вашего текста указано, что потенциал $u$ вводится как \textbf{интенсивность субстанции}, создаваемой притягивающим телом. Математически ввод потенциала строится на следующих этапах:

    \begin{enumerate}
        \item \textbf{Определение для точечного источника:} Для одной массы $m$ в точке $\mathbf{x}_0$ потенциал в произвольной точке $\mathbf{x}$ вводится формулой (3.1):
        \begin{equation}
        u(\mathbf{x}) = \gamma \frac{m}{|\mathbf{x} - \mathbf{x}_0|}
        \end{equation}
        
        \item \textbf{Связь с силой (Градиент):} Потенциал вводится таким образом, чтобы вектор силы притяжения $\mathbf{f}$ можно было найти как его градиент (формула 3.3):
        \begin{equation}
        \mathbf{f} = \text{grad} u
        \end{equation}
        Это означает, что компоненты силы — это частные производные потенциала по координатам: $f_x = \frac{\partial u}{\partial x}$, $f_y = \frac{\partial u}{\partial y}$, $f_z = \frac{\partial u}{\partial z}$.

        \item \textbf{Принцип суперпозиции:} Если источников много, потенциалы от каждого из них просто складываются (формула 3.5):
        \begin{equation}
        u(\mathbf{x}) = \gamma \sum_{j=1}^{N} \frac{m_j}{|\mathbf{x} - \mathbf{x}_j|}
        \end{equation}

        \item \textbf{Непрерывное распределение (Объемный потенциал):} При переходе от суммы к пределу для массы, распределенной с плотностью $\rho$, вводится интегральное определение (формула 3.11):
        \begin{equation}
        u(\mathbf{x}_0) = \gamma \int_{\Omega} \frac{\rho(\mathbf{x}) d\mathbf{x}}{|\mathbf{x} - \mathbf{x}_0|}
        \end{equation}
    \end{enumerate}
\end{itemize}
  \item Выведите уравнение Лапласа для потенциала точечной массы.\section*{ГИПОТЕЗА ДАЛЬНОДЕЙСТВИЯ ЛАПЛАСА}

Наличие любого притягивающего (т. е. с положительной массой) тела влечет за собой возникновение во всем пространстве некоторой субстанции, интенсивность $u$ которой в произвольной точке $\mathbf{x} = (x, y, z)$ вычисляется по формуле:

\begin{equation}
u(\mathbf{x}) = \gamma \frac{m}{|\mathbf{x} - \mathbf{x}_0|}
\end{equation}

Здесь $m$ — масса данного тела, $\mathbf{x}_0 = (x_0, y_0, z_0)$ — место (точка) расположения тела, $\gamma$ — абсолютная константа, получившая название гравитационной постоянной, ее численное значение в системе СИ равно $6.6732 \cdot 10^{-11} \text{Н} \cdot \text{м}^2 / \text{кг}^2$, наконец, $|\mathbf{x} - \mathbf{x}_0|$ — расстояние между точками $\mathbf{x}$ и $\mathbf{x}_0$, определяемое формулой:

\begin{equation}
|\mathbf{x} - \mathbf{x}_0| = \sqrt{(x - x_0)^2 + (y - y_0)^2 + (z - z_0)^2}
\end{equation}

Смысл субстанции $u$ заключается в том, что ее знание позволяет вычислить вектор $\mathbf{f} = (f_x, f_y, f_z)$ силы притяжения, действующей со стороны тела на тело единичной массы, расположенное в точке $\mathbf{x}$, с помощью формулы:

\begin{equation}
\mathbf{f} = \text{grad} u
\end{equation}

Вектор $\text{grad} u$ в (3.3) называется \textit{градиентом} функции $u$. По своему смыслу $\text{grad} u(\mathbf{x})$ указывает направление быстрейшего возрастания функции $u$ в точке $\mathbf{x}$. В декартовой системе координат с ортами $\mathbf{i}, \mathbf{j}, \mathbf{k}$ он определяется формулой:

\begin{equation}
\text{grad} u = \frac{\partial u}{\partial x} \mathbf{i} + \frac{\partial u}{\partial y} \mathbf{j} + \frac{\partial u}{\partial z} \mathbf{k}
\end{equation}

С учетом это векторное равенство (3.3) можно записать в эквивалентном виде трех скалярных равенств для декартовых компонент $f_x, f_y$ и $f_z$:

\begin{equation}
f_x = \frac{\partial u}{\partial x}, \quad f_y = \frac{\partial u}{\partial y}, \quad f_z = \frac{\partial u}{\partial z}
\end{equation}

Поскольку вектор $\mathbf{f}$ описывает силу тяготения, то функцию $u$ принято называть потенциалом поля силы тяготения или просто \textit{гравитационным потенциалом}. При этом само поле часто называют гравитационным полем.

Подчеркнем, что для гравитационного поля справедлив принцип суперпозиции. Согласно этому принципу гравитационное поле, создаваемое совокупностью притягивающих тел, равно сумме гравитационных полей, создаваемых каждым из этих тел. Отсюда вытекает, что гравитационный потенциал поля, создаваемого $N$ парами $(\mathbf{x}_1, m_1), (\mathbf{x}_2, m_2), \dots, (\mathbf{x}_N, m_N)$, определяется в произвольной точке $\mathbf{x} \neq \mathbf{x}_j$ формулой:

\begin{equation}
u(\mathbf{x}) = \gamma \sum_{j=1}^{N} \frac{m_j}{|\mathbf{x} - \mathbf{x}_j|}
\end{equation}

---

\textit{Вывод уравнения Лапласа (стр. 31):}

Выберем $i$-ое слагаемое в сумме (3.5) и вычислим вторые производные от функции $u_i$. Полагая для простоты $|\mathbf{x} - \mathbf{x}_i| = r_i(\mathbf{x})$, имеем:

\begin{equation}
\frac{\partial r_i(\mathbf{x})}{\partial x} = \frac{x - x_i}{r}, \quad \frac{\partial r_i(\mathbf{x})}{\partial y} = \frac{y - y_i}{r}, \quad \frac{\partial r_i(\mathbf{x})}{\partial z} = \frac{z - z_i}{r}
\end{equation}

Отсюда и из (3.6) вытекает, что:

\begin{equation}
\frac{\partial u_i(\mathbf{x})}{\partial x} = -\gamma m_i \frac{x - x_i}{r^3}, \quad \frac{\partial u_i(\mathbf{x})}{\partial y} = -\gamma m_i \frac{y - y_i}{r^3}, \quad \frac{\partial u_i(\mathbf{x})}{\partial z} = -\gamma m_i \frac{z - z_i}{r^3}
\end{equation}

Дифференцируя еще раз, имеем:

\begin{equation}
\frac{\partial^2 u_i(\mathbf{x})}{\partial x^2} = \gamma m_i \left[ -\frac{1}{r^3} + 3 \frac{(x - x_i)^2}{r^5} \right], \quad \dots \text{ (аналогично для y и z)}
\end{equation}

Складывая найденные частные производные, получаем:

\begin{equation}
\frac{\partial^2 u_i}{\partial x^2} + \frac{\partial^2 u_i}{\partial y^2} + \frac{\partial^2 u_i}{\partial z^2} = 0, \quad \mathbf{x} \neq \mathbf{x}_i
\end{equation}

Отсюда и из того условия, что $u = \sum_{i=1}^N u_i$, приходим к уравнению:

\begin{equation}
\Delta u = \frac{\partial^2 u}{\partial x^2} + \frac{\partial^2 u}{\partial y^2} + \frac{\partial^2 u}{\partial z^2} = 0
\end{equation}
  \item Что такое объемный потенциал? Запишите уравнение Пуассона для гравитационного и электростатического полей.\begin{itemize}
    \item \textbf{Что такое объемный потенциал?}
    
    На практике часто приходится иметь дело с полем тяготения, вызванным массой, распределенной с плотностью $\rho$ по некоторой области $\Omega$. Если разбить область $\Omega$ на элементарные подобласти $\Omega_i$ с объемами $\Delta V_i$ и применить принцип суперпозиции, то потенциал в произвольной точке $\mathbf{x}_0$ будет определяться интегралом.
    
    \textbf{Объемным (или ньютоновским) потенциалом} называется функция, определяемая формулой:
    \begin{equation}
    u(\mathbf{x}_0) = \gamma \int_{\Omega} \frac{\rho(\mathbf{x}) d\mathbf{x}}{|\mathbf{x} - \mathbf{x}_0|} \tag{3.11}
    \end{equation}
    где $d\mathbf{x}$ — элемент объема, а $\gamma$ — гравитационная постоянная. Этот потенциал точно определяет поле тела $(\Omega, \rho)$ в точках $\mathbf{x}_0$, расположенных вне области $\Omega$.

    \item \textbf{Уравнение Пуассона для гравитационного и электростатического полей.}
    
    Для плотности $\rho$, обладающей достаточной гладкостью, потенциал $u$ удовлетворяет в каждой точке $\mathbf{x}$ пространства $\mathbb{R}^3$ \textbf{уравнению Пуассона}. Оно является уравнением эллиптического типа.
    
    \textbf{1. Для гравитационного поля:}
    \begin{equation}
    \Delta u \equiv \frac{\partial^2 u}{\partial x^2} + \frac{\partial^2 u}{\partial y^2} + \frac{\partial^2 u}{\partial z^2} = -4\pi\gamma\rho(\mathbf{x}) \tag{3.12}
    \end{equation}
    Здесь $\gamma$ — гравитационная постоянная, $\rho$ — плотность распределения массы. Вне притягивающих масс (где $\rho = 0$) данное уравнение переходит в уравнение Лапласа ($\Delta u = 0$).

    \textbf{2. Для электростатического поля (по аналогии):}
    Уравнение Пуассона связывает электростатический потенциал $\phi$ с плотностью электрического заряда $\rho_e$:
    \begin{equation}
    \Delta \phi = -\frac{\rho_e}{\varepsilon_0}
    \end{equation}
    где $\varepsilon_0$ — электрическая постоянная.
\end{itemize}
  \end{enumerate}

\item \textbf{Граничные условия}
  \begin{enumerate}[label=(\alph*)]
  \item Сформулируйте условия Дирихле, Неймана и третье краевое условие для уравнения Пуассона. Приведите физические примеры.\begin{itemize}
    \item \textbf{Граничные условия для уравнения Пуассона}
    
    Для выделения единственного решения уравнения Пуассона в области $\Omega$ необходимо задать дополнительные условия на ее границе $\Gamma$. Согласно тексту, выделяют три основных типа условий:

    \begin{enumerate}
        \item \textbf{Условие 1-го рода (Условие Дирихле):}
        Задается значение самого потенциала $u$ на границе:
        \begin{equation}
        u = g \quad \text{на } \Gamma \tag{3.20}
        \end{equation}
        \textit{Физический пример:} Если $g=0$, это означает, что граница $\Gamma$ является \textbf{заземленной} или эквипотенциальной (служит экраном для электрического поля).

        \item \textbf{Условие 2-го рода (Условие Неймана):}
        Задается значение нормальной производной потенциала на границе:
        \begin{equation}
        \frac{\partial u}{\partial n} = g \quad \text{на } \Gamma \tag{3.21}
        \end{equation}
        \textit{Физический пример:} Поскольку градиент потенциала связан с напряженностью поля, это условие соответствует заданию \textbf{потока} электрической индукции или плотности заряда на границе проводника.

        \item \textbf{Условие 3-го рода (Третье краевое условие):}
        Представляет собой линейную комбинацию потенциала и его нормальной производной:
        \begin{equation}
        \frac{\partial u}{\partial n} + \alpha u = g \quad \text{на } \Gamma \tag{3.22}
        \end{equation}
    \end{enumerate}

    \item \textbf{Общее граничное условие}
    
    Все вышеперечисленные случаи объединяются общим уравнением (3.23):
    \begin{equation}
    \alpha u + \beta \frac{\partial u}{\partial n} = g \quad \text{на } \Gamma \tag{3.23}
    \end{equation}
    Где при $\alpha=1, \beta=0$ получаем задачу Дирихле, а при $\alpha=0, \beta=1$ — задачу Неймана.

    \item \textbf{Важное замечание о гравитации (Замечание 3.2):}
    В тексте подчеркивается, что данные граничные условия имеют физический смысл в основном для \textbf{электростатического} потенциала. Для гравитационного потенциала уравнение Пуассона рассматривается во всем пространстве $\mathbb{R}^3$ без граничных условий, так как невозможно создать экран, полностью изолирующий область от действия гравитации.
\end{itemize}
  \item Почему для гравитационного поля обычно не ставят граничных условий на конечной границе?
  \end{enumerate}

\item \textbf{Уравнение теплопроводности}
  \begin{enumerate}[label=(\alph*)]
  \item Запишите закон Фурье. Какой физический смысл имеет вектор потока тепла?
  \item Выведите уравнение баланса тепла, используя формулу Гаусса--Остроградского. 




ности будем рассматривать лишь жидкие среды.

Предположим, что жидкость, находящаяся в состоянии покоя, занимает некоторую область $D$ пространства $\mathbb{R}^3$. Обозначим через $\Omega$ произвольную ее ограниченную подобласть с кусочно-гладкой границей $\Gamma$. Для математического описания процесса переноса тепла принято вводить \textit{вектор потока тепла} $\mathbf{q} = \mathbf{q}(\mathbf{x}, t)$. Физический смысл вектора $\mathbf{q}$ заключается в том, что с его помощью можно определить количество $Q_1$ тепла, вносимое за промежуток времени от $t_1$ до $t_2$ в произвольную подобласть $\Omega$ рассматриваемой области $D$ со стороны оставшейся части $\Omega_e = D \setminus \overline{\Omega}$, по формуле:
\begin{equation}\tag{4.1}
Q_1 = -\int_{t_1}^{t_2} dt \int_{\Gamma} \mathbf{q} \cdot \mathbf{n} \, dS.
\end{equation}
Здесь $\Gamma$ --- граница между $\Omega$ и $\Omega_e$ (рис. 4.1), $\mathbf{n}$ --- единичный вектор внешней нормали к поверхности $\Gamma$, $dS$ --- элемент площади поверхности $\Gamma$. Знак <<$-$>> выбран с учетом ориентации нормали $\mathbf{n}$.

Из школьного курса термодинамики хорошо известно, что мерой тепла является температура $T$. В связи с этим возникает принципиально важный вопрос о том, как связать поток тепла с температурой $T$. Основополагающую роль в термодинамике играет закон Фурье, названный так в честь известного французского математика и физика S.B. Fourier (1768--1830).
Указанный закон постулирует, что вектор потока тепла связан с температурой $T$ формулой
\begin{equation}\tag{4.2}
\mathbf{q} = -k \mathrm{grad} T.
\end{equation}
Здесь $k$ --- параметр состояния среды, называемый \textit{коэффициентом теплопроводности}, а знак <<$-$>> в законе Фурье отражает тот известный опытный факт, что тепло всегда течет от горячей части среды к холодной. В словесной форме закон Фурье выражается следующими словами: \textit{если температура тела не равномерна, то в нем возникают тепловые потоки, направленные из мест с более высокой температурой в места с более низкой температурой}. Размерности величин $\mathbf{q}$, $k$ и $T$ в системе СИ приведены в таблице 4.1.

Напомним, что занимающая область $D$ среда называется \textit{однородной}, если ее свойства не меняются при переходе от одной точки $\mathbf{x} \in D$ к другой. Среда называется \textit{изотропной} в точке $\mathbf{x}$, если ее свойства одинаковы по всем направлениям, выходящим из точки $\mathbf{x}$, и \textit{анизотропной} в противном случае. Если среда изотропна, то $k$ является скалярной величиной, зависящей в общем случае от точек $\mathbf{x} \in D$, времени $t$ и температуры $T$. В случае анизотропной среды $k$ является тензорной функцией, а вектор потока тепла представляет собой (скалярное) произведение тензора $k$ на вектор $-\mathrm{grad} T$. Ниже мы будем рассматривать только изотропные среды. Для таких сред формула (4.1) принимает с учетом (4.2) вид
\[
Q_1 = \int_{t_1}^{t_2} dt \int_{\Gamma} k \mathrm{grad} T \cdot \mathbf{n} \, dS.
\]
Используя формулу Гаусса-Остроградского [19, с.192]
\begin{equation}\tag{4.3}
\int_{\Omega} \mathrm{div} \mathbf{v} \, d\mathbf{x} = \int_{\Gamma} \mathbf{v} \cdot \mathbf{n} \, dS
\end{equation}
для дифференцируемого в $\Omega$ векторного поля $\mathbf{v}$, в которой следует положить $\mathbf{v} = k \mathrm{grad} T$, выражение для $Q_1$ можно переписать в виде
\begin{equation}\tag{4.4}
Q_1 = \int_{t_1}^{t_2} dt \int_{\Omega} \mathrm{div}(k \mathrm{grad} T) \, d\mathbf{x}.
\end{equation}
Формула (4.4) описывает количество тепла, поступающее в область $\Omega$ за промежуток времени от $t_1$ до $t_2$, вследствие потоков тепла, вызываемых неравномерным распределением температуры в $D$.

Предположим далее, что в области $\Omega$ распределены источники тепла с объемной плотностью $F$. Тогда количество тепла $Q_2$, выделяемое ими в область $\Omega$ за время $\Delta t = t_2 - t_1$, будет равно
\[
Q_2 = \int_{t_1}^{t_2} dt \int_{\Omega} F \, d\mathbf{x}.
\]
Ясно, что получаемое областью $\Omega$ тепло идет на нагрев среды, т. е. на увеличение ее температуры. Если в момент $t_1$ температура в произвольной точке $\mathbf{x}$ области $\Omega$ равна $T_1(\mathbf{x})$, а в момент $t_2$ равна $T_2(\mathbf{x})$, то количество тепла, необходимое для указанного увеличения температуры, равно в силу законов термодинамики $Q$, где
\begin{equation}\tag{4.5}
Q = \int_{\Omega} \rho c (T_2 - T_1) \, d\mathbf{x} = \int_{\Omega} \rho c \int_{t_1}^{t_2} \frac{\partial T}{\partial t} \, dt \, d\mathbf{x} = \int_{t_1}^{t_2} dt \int_{\Omega} \rho c \frac{\partial T}{\partial t} \, d\mathbf{x}.
\end{equation}
Здесь параметры $\rho$ (плотность среды) и $c$ (коэффициент удельной теплоемкости при постоянном давлении) являются в общем случае функциями от $\mathbf{x} \in D$.

В силу \textit{фундаментального закона сохранения тепла}, примененного к области $\Omega$, должно выполняться соотношение: $Q = Q_1 + Q_2$, или
\begin{equation}\tag{4.6}
\int_{t_1}^{t_2} dt \int_{\Omega} \rho c \frac{\partial T}{\partial t} \, d\mathbf{x} = \int_{t_1}^{t_2} dt \int_{\Omega} \mathrm{div}(k \mathrm{grad} T) \, d\mathbf{x} + \int_{t_1}^{t_2} dt \int_{\Omega} F \, d\mathbf{x}.
\end{equation}
Оно называется \textit{уравнением баланса тепла}. В словесной форме соотношение (4.6) можно выразить так: изменение количества тепла в области $\Omega$ за время $\Delta t = t_2 - t_1$ (см. левую часть в (4.6)) обусловлено притоком тепла в область $\Omega$ через граничную поверхность $\Gamma$ за счет молекулярной диффузии, и количеством тепла, выделившимся в $\Omega$ за время $\Delta t$ в результате действия объемных источников тепла с плотностью $F$.

Переписав соотношение (4.6) в виде
\[
\int_{t_1}^{t_2} dt \int_{\Omega} \left[ \rho c \frac{\partial T}{\partial t} - \mathrm{div}(k \mathrm{grad} T) - F \right] d\mathbf{x} = 0,
\]
предположим, что подынтегральная функция здесь является непрерывной в области $D$. Поскольку в этом соотношении $\Omega$ --- произвольная область с кусочно-гладкой границей $\Gamma$, то тем самым мы находимся в условиях леммы 1.1. Согласно этой лемме указанная подынтегральная функция равна нулю всюду в области $D$, т. е. выполняется соотношение
\begin{equation}\tag{4.7}
\rho c \frac{\partial T}{\partial t} = \mathrm{div}(k \mathrm{grad} T) + F.
\end{equation}
Уравнение (4.7) и является искомой математической моделью, описывающей процесс распространения тепла в области $D$. Как уже указывалось, коэффициенты $\rho$ и $c$ в (4.7) являются в общем случае функциями от $\mathbf{x}$, а коэффициент $k$ к тому же может зависеть и от $t$, а также самой температуры $T$. Поэтому (4.7) является в общем случае нелинейным уравнением.


  \item Получите дифференциальное уравнение теплопроводности. Что такое коэффициент температуропроводности?
  \item Какие начальные и граничные условия ставятся для уравнения теплопроводности?В декартовой системе координат уравнение (4.7) имеет вид
\begin{equation}\tag{4.8}
\rho c \frac{\partial T}{\partial t} = \frac{\partial}{\partial x} \left( k \frac{\partial T}{\partial x} \right) + \frac{\partial}{\partial y} \left( k \frac{\partial T}{\partial y} \right) + \frac{\partial}{\partial z} \left( k \frac{\partial T}{\partial z} \right) + F.
\end{equation}

В частном случае, когда среда однородна, так что величины $\rho, c$ и $k$ есть константы, уравнение (4.7) принимает вид
\begin{equation}\tag{4.9}
\frac{\partial T}{\partial t} = a^2 \Delta T + f,
\end{equation}
где $a = \sqrt{k/\rho c}$, $f = F/\rho c$. Уравнение (4.9) называется \textit{неоднородным уравнением теплопроводности}, а $a^2$ называется \textit{коэффициентом температуропроводности}. Если объемные источники тепла отсутствуют, так что $f = 0$, то уравнение (4.9) становится однородным. Указанное уравнение
\begin{equation}\tag{4.10}
\frac{\partial T}{\partial t} = a^2 \Delta T
\end{equation}
является важнейшим и в то же время простейшим представителем уравнений параболического типа согласно классификации, приведенной в гл. 2.

Для выделения единственного решения уравнения (4.7) необходимо задавать дополнительные условия. В случае, когда процесс распространения тепла рассматривается во всем пространстве $\mathbb{R}^3$, так что $\Omega = \mathbb{R}^3$, роль таких условий играет начальное условие
\begin{equation}\tag{4.11}
T|_{t=0} = T_0(\mathbf{x}), \quad \mathbf{x} \in \Omega,
\end{equation}
а задача (4.7), (4.11) при $\Omega = \mathbb{R}^3$ называется \textit{задачей Коши} для уравнения (4.7). Если процесс распространения тепла рассматривается в некоторой области $\Omega$ пространства $\mathbb{R}^3$ с границей $\Gamma = \partial\Omega$, то на границе $\Gamma$ следует задавать краевые условия, имеющие в общем случае вид (ср. с (3.23)):
\begin{equation}\tag{4.12}
\alpha T + \beta \frac{\partial T}{\partial n} = g \quad \text{на } \Gamma.
\end{equation}
Здесь $\alpha, \beta$ и $g$ --- заданные на $\Gamma$ функции. Задача (4.7), (4.11), (4.12) называется \textit{начально-краевой задачей} для уравнения (4.7): \textit{первой краевой задачей} или \textit{задачей Дирихле} при $\beta = 0, \alpha = 1$, \textit{второй краевой задачей} или \textit{задачей Неймана} при $\beta = k|_{\Gamma}, \alpha = 0$ и \textit{третьей краевой задачей} при $\beta = k|_{\Gamma}, \alpha \neq 0$.

В физическом плане условие Дирихле ($\alpha = 1, \beta = 0$) в (4.12) отвечает заданию температуры на границе $\Gamma$, условие Неймана ($\alpha = 0, \beta = k|_{\Gamma}$) в (4.12) отвечает заданию теплового потока на $\Gamma$, наконец, условие 3-го рода означает, что на границе происходит теплообмен с внешней средой, температура которой известна. Подробнее о физическом смысле различных граничных условий для температуры $T$ можно прочитать в [21, с.27--28], [22, с.15--16], [56, с.196--202].

В частном случае, когда температура $T$ зависит лишь от координат $x, y$ и времени $t$, что имеет место, например, при изучении распространения тепла в тонкой однородной пластине или мембране, (4.9) принимает вид \textit{двумерного уравнения теплопроводности}
\begin{equation}\tag{4.13}
\frac{\partial T}{\partial t} = a^2 \left( \frac{\partial^2 T}{\partial x^2} + \frac{\partial^2 T}{\partial y^2} \right) + f.
\end{equation}

Наконец, для одномерного однородного тела, например, для однородного стержня, когда $T$ зависит только от $x$ и $t$, (4.9) переходит в \textit{одномерное уравнение теплопроводности}
\begin{equation}\tag{4.14}
\frac{\partial T}{\partial t} = a^2 \frac{\partial^2 T}{\partial x^2} + f.
\end{equation}

Следует, однако, отметить, что уравнения (4.13) и (4.14) не учитывают теплообмен между поверхностью пластинки или стержня с окружающей средой (точнее мы пренебрегаем им при выводе этих уравнений).

В ряде случаев зависимостью температуры $T$ от времени $t$ можно пренебречь. В этих случаях (4.7) либо (4.8) переходит в стационарное уравнение
\begin{equation}\tag{4.15}
\mathrm{div}(k \mathrm{grad} T) \equiv \frac{\partial}{\partial x} \left( k \frac{\partial T}{\partial x} \right) + \frac{\partial}{\partial y} \left( k \frac{\partial T}{\partial y} \right) + \frac{\partial}{\partial z} \left( k \frac{\partial T}{\partial z} \right) = -F.
\end{equation}

В случае, когда $k = \text{const}$, (4.15) переходит в уравнение Пуассона
\begin{equation}\tag{4.16}
\Delta T = -f,
\end{equation}
где $f = F/k$. В обоих случаях начальное условие (4.11), естественно, снимается, так что изучение процесса переноса тепла сводится к решению краевой задачи для уравнения с переменными коэффициентами (4.15) либо уравнения Пуассона (4.16). В простейшем случае, отвечающем распределению температуры в одномерном стержне длины $l$, уравнение (4.16) вместе с простейшим граничным условием 1-го рода в (4.12) принимает вид краевой задачи для обыкновенного дифференциального уравнения 2-го порядка:
\begin{equation}\tag{4.17}
\frac{d^2 T}{dx^2} = -f, \quad T(0) = T(l) = 0.
\end{equation}
Граничные условия в (4.17) означают, что температура $T$ на концах стержня, в котором изучается тепловой процесс, поддерживается равной нулю.

Уравнения (4.7), (4.15) относятся к случаю, когда жидкая среда находится в покое. Если жидкость, занимающая область $\Omega$, находится в движении, то возникает дополнительный механизм переноса тепла за счет движущихся частиц жидкости. В физике указанный механизм называется \textit{(тепловой) конвекцией}. Для вывода соответствующей математической модели...
  \end{enumerate}

\item \textbf{Модели движения жидкости}
  \begin{enumerate}[label=(\alph*)]
  \item Что такое модель идеальной жидкости? Запишите уравнение неразрывности и уравнение Эйлера.Считая подынтегральную функцию в (5.3) непрерывной и применяя лемму 1.1, приходим к дифференциальному уравнению
\begin{equation}\tag{5.4}
\frac{\partial \psi}{\partial t} + \mathrm{div}(\psi \mathbf{u}) = q.
\end{equation}
Уравнение (5.4) называется \textit{дифференциальным законом сохранения} величины $\Psi$, тогда как уравнение (5.1) имеет смысл \textit{интегрального закона сохранения} величины $\Psi$. Применяя (5.4) к конкретным гидродинамическим величинам, теперь нетрудно вывести законы сохранения массы, импульса и т. д., на основе которых выводятся искомые математические модели движения жидкости.

\subsection*{5.3. Модели движения идеальной жидкости}
Начнем с вывода закона сохранения массы. С этой целью положим в (5.4) $\psi = \rho, q = 0$ (последнее объясняется отсутствием внешних источников массы). В результате получим уравнение
\begin{equation}\tag{5.5}
\frac{\partial \rho}{\partial t} + \mathrm{div}(\rho \mathbf{u}) = 0.
\end{equation}
Оно называется \textit{дифференциальным законом сохранения массы}, или просто \textit{уравнением неразрывности}. В силу закона сохранения массы движение жидкости происходит так, что ее скорость $\mathbf{u}$ и плотность $\rho$ в каждой точке $\mathbf{x} \in D$ в любой момент $t$ удовлетворяют уравнению (5.5). Важно отметить, что (5.5) имеет инвариантный характер. Однако для упрощения дальнейших выкладок нам будет удобно считать, что в области $D$ введена декартова система координат $x_1 = x, x_2 = y, x_3 = z$ с правым базисом $\mathbf{i, j, k}$. Обозначив декартовы компоненты вектора $\mathbf{u}$ через $u_1 = u, u_2 = v, u_3 = w$ и учитывая формулу (3.16) для дивергенции, уравнение (5.5) можно переписать в виде
\begin{equation}\tag{5.6}
\frac{\partial \rho}{\partial t} + \sum_{i=1}^3 \frac{\partial}{\partial x_i} (\rho u_i) = 0.
\end{equation}

Прежде чем приступить к выводу уравнения сохранения импульса, рассмотрим силы, действующие на жидкость. В общем случае эти силы разделяются на \textit{внешние} и \textit{внутренние}. Внешние силы, как правило, являются \textit{массовыми} (либо \textit{объемными}) и описываются массовой плотностью $\mathbf{f}$. Таким образом, если $\mathbf{f}$ --- массовая плотность внешней силы, то $\rho \mathbf{f}$ представляет собой ее \textit{объемную плотность}. $\rho \mathbf{f} d\mathbf{x}$ --- силу, действующую на элементарный объем $d\mathbf{x}$. Важным примером внешних сил является сила тяжести:
\begin{equation}\tag{5.7}
\mathbf{g} = (0, 0, -g).
\end{equation}

Обратившись к анализу внутренних сил, заметим, что они существенно зависят от характера жидкости. Мы будем рассматривать \textit{идеальную жидкость}, т. е. такую среду, в которой внутренние силы сводятся к давлению $p$. Сила, действующая на элемент поверхности $dS$ в идеальной жидкости:
\begin{equation}\tag{5.8}
-p \mathbf{n} dS = -(p \cos \alpha \mathbf{i} + p \cos \beta \mathbf{j} + p \cos \gamma \mathbf{k}) dS,
\end{equation}
где $\cos \alpha, \cos \beta, \cos \gamma$ --- направляющие косинусы внешней нормали $\mathbf{n}$. На всю область $\Omega$ будет действовать сила, равная поверхностному интегралу
\begin{equation}\tag{5.9}
-\int_{\Gamma} p \mathbf{n} dS = -\int_{\Omega} \nabla p d\mathbf{x}.
\end{equation}
Из (5.9) следует, что объемная плотность сил в идеальной жидкости имеет вид
\begin{equation}\tag{5.10}
-\nabla p + \rho \mathbf{f}.
\end{equation}

В проекциях на оси координат (5.10) имеет вид
\begin{equation}\tag{5.11}
-\frac{\partial p}{\partial x_i} + \rho f_i, \quad i = 1, 2, 3.
\end{equation}
В результате получим уравнение
\begin{equation}\tag{5.12}
\frac{\partial}{\partial t}(\rho u_i) + \mathrm{div}(\rho u_i \mathbf{u}) = -\frac{\partial p}{\partial x_i} + \rho f_i.
\end{equation}
В силу уравнения неразрывности (5.6) выражение в квадратных скобках обращается в нуль, и мы приходим к \textit{уравнениям Эйлера}:
\begin{equation}\tag{5.13}
\rho \frac{\partial u_i}{\partial t} + \rho \sum_{j=1}^3 u_j \frac{\partial u_i}{\partial x_j} = -\frac{\partial p}{\partial x_i} + \rho f_i, \quad i = 1, 2, 3
\end{equation}
или в векторной форме:
\begin{equation}\tag{5.14}
\rho \frac{\partial \mathbf{u}}{\partial t} + \rho (\mathbf{u} \cdot \nabla) \mathbf{u} = -\nabla p + \rho \mathbf{f}.
\end{equation}

Для замыкания системы используется уравнение состояния. Например, для баротропной жидкости:
\begin{equation}\tag{5.15}
p = P(\rho).
\end{equation}
Более реалистичными являются двухпараметрические среды $p = P(\rho, s)$, где $s$ --- энтропия. Для них добавляется уравнение:
\begin{equation}\tag{5.17}
\frac{\partial s}{\partial t} + \mathrm{div}(s \mathbf{u}) = 0.
\end{equation}

\subsection*{5.4. Модели движения вязкой жидкости}
Для вязкой жидкости необходимо учитывать силы трения. В простейшем случае объемная плотность сил вязкого трения моделируется выражением $\mu \Delta \mathbf{u}$. Добавив его в уравнение движения (5.14), приходим к \textit{уравнению Навье-Стокса}:
\begin{equation}\tag{5.23}
\rho \frac{\partial \mathbf{u}}{\partial t} + \rho (\mathbf{u} \cdot \nabla) \mathbf{u} = -\nabla p + \mu \Delta \mathbf{u} + \rho \mathbf{f}.
\end{equation}
  \item Как выводится уравнение Навье--Стокса? В чем состоит шестая проблема тысячелетия?\subsection*{5.4. Модели движения вязкой жидкости}

Модели идеальной жидкости являются весьма приближенными, поскольку в реальных жидкостях всегда присутствует трение, вызываемое наличием вязкости в жидкости. Наличие вязкости в жидкости приводит к возникновению дополнительной внутренней силы (препятствующей движению жидкости). Ее объемная плотность часто моделируется выражением [40]:
\begin{equation}\tag{5.22}
\mu \Delta \mathbf{u}.
\end{equation}
Здесь $\Delta \mathbf{u}$ --- векторный лапласиан от скорости $\mathbf{u}$, определяемый, в частности, в декартовой системе координат формулой $\Delta \mathbf{u} = \Delta u \mathbf{i} + \Delta v \mathbf{j} + \Delta w \mathbf{k}$, $\mu$ --- постоянный коэффициент динамической вязкости, имеющий размерность кг/м$\cdot$сек. В общем случае коэффициент $\mu$ может зависеть как от точек $\mathbf{x} \in D$, так и от некоторых характеристик среды, например, температуры, причем выражение для силы вязкого трения имеет более сложный вид, чем в (5.22) (см., например, [40]). Однако рассмотрение более сложных моделей выходит за рамки данной книги. Добавив выражение (5.22) в правую часть основного уравнения движения идеальной жидкости (5.14), приходим к уравнению
\begin{equation}\tag{5.23}
\rho \frac{\partial \mathbf{u}}{\partial t} + \rho (\mathbf{u} \cdot \nabla) \mathbf{u} = -\nabla p + \mu \Delta \mathbf{u} + \rho \mathbf{f},
\end{equation}
представляющему собой основное уравнение движения вязкой жидкости. Оно называется \textit{векторным уравнением Навье--Стокса} в честь французского инженера A. Navier (1785--1836) и английского физика G.G. Stokes (1819--1903), много сделавших для становления и развития гидродинамики вязкой жидкости.

Заменив в моделях $M_1$ и $M_5$ идеальных жидкостей уравнение движения идеальной жидкости (5.14) уравнением (5.23), мы получим еще две математические модели, описывающие движение вязкой жидкости: модель $M_6$ вязкой баротропной жидкости и модель $M_7$ вязкой несжимаемой однородной жидкости. Они описываются соответственно уравнениями
\[
\rho \frac{\partial \mathbf{u}}{\partial t} + \rho (\mathbf{u} \cdot \nabla) \mathbf{u} = -\nabla p + \mu \Delta \mathbf{u} + \rho \mathbf{f}, \quad \frac{\partial \rho}{\partial t} + \text{div}(\rho \mathbf{u}) = 0, \quad p = P(\rho),
\]
и
\[
\rho \frac{\partial \mathbf{u}}{\partial t} + \rho (\mathbf{u} \cdot \nabla) \mathbf{u} = -\nabla p + \mu \Delta \mathbf{u} + \rho \mathbf{f}, \quad \text{div} \mathbf{u} = 0, \quad \rho = \text{const}.
\]
Последняя модель, называемая \textit{моделью Навье--Стокса}, является одной из важнейших моделей в гидродинамике вязкой несжимаемой жидкости.

Наряду с приведенными моделями движения вязкой жидкости, существует еще большое количество других моделей, в которых, кроме сил вязкого трения, учитываются и другие эффекты, описываемые дифференциальными операторами второго порядка. Среди них различают модели теплопереноса, учитывающие действие в жидкости как сил трения, так и тепловых эффектов, модели массопереноса, модели вязкой токопроводящей жидкости, учитывающие эффекты действия электромагнитных сил на проводящую жидкость, и т. д.

Ограничимся здесь приведением трех математических моделей. Первая модель описывает движение вязкой однородной несжимаемой теплопроводной жидкости. Предварительно отметим, что наличие градиента температуры в жидкости приводит к возникновению в ней дополнительной объемной силы. Чтобы вывести выражение для ее плотности, рассмотрим область $D$, в которой находится жидкость с постоянной плотностью $\rho_0$ и постоянной температурой $T_0$. Предположим, что после нагревания либо охлаждения некоторой подобласти $\Omega$ области $D$ плотность и температура жидкости в $\Omega$ становятся равными $\rho$ и $T$. Тогда на жидкость в области $\Omega$ со стороны оставшейся жидкости будет действовать сила, имеющая смысл силы плавучести, объемная плотность которой равна $(\rho - \rho_0)\mathbf{g}$. Если разность $\rho - \rho_0$ мала по сравнению с $\rho_0$, то термодинамическое уравнение состояния жидкости может быть приближенно записано в виде $\rho = \rho_0 - \rho_0 \alpha_T(T - T_0)$ [5, с.210], [15, с.264]. Здесь параметр $\alpha_T > 0$, называемый \textit{объемным коэффициентом теплового расширения}, имеет размерность $K^{-1}$, где $K$ --- сокращенное обозначение Кельвина как единицы температуры. С учетом указанного приближения объемную плотность $\mathbf{f}_T$ силы плавучести можно записать в виде
\begin{equation}\tag{5.24}
\mathbf{f}_T = -\rho_0 \alpha_T(T - T_0)\mathbf{g}.
\end{equation}

Выражение (5.24) следует добавить в уравнение (5.23), в котором следует положить $\rho = \rho_0$. Поскольку в (5.24) входит неизвестная в общем случае температура $T$, то для построения замкнутой модели к построенной системе уравнений необходимо присоединить соответствующее уравнение для $T$. В результате приходим к следующей модели
\[
\rho_0 \frac{\partial \mathbf{u}}{\partial t} + \rho_0 (\mathbf{u} \cdot \nabla) \mathbf{u} = -\nabla p + \mu \Delta \mathbf{u} + \rho_0 \mathbf{f} - \rho_0 \alpha_T(T - T_0)\mathbf{g}, \quad \text{div} \mathbf{u} = 0,
\]
  \item Что такое безвихревое течение? Каким уравнением удовлетворяет потенциал скорости?\subsection*{5.5. Модель стационарного безвыхревого движения несжимаемой жидкости}

Рассмотрим стационарный процесс движения идеальной несжимаемой однородной жидкости под действием силы тяжести. Отвечающая этому случаю модель $M_5$ принимает с учетом замечания 5.1 и условия $\mathbf{f} = \mathbf{g}$ вид:
\begin{equation}\tag{5.29}
(\mathbf{u} \cdot \nabla)\mathbf{u} = -\frac{1}{\rho} \nabla p + \mathbf{g}, \quad \text{div} \mathbf{u} = 0, \quad \rho = \text{const}.
\end{equation}


Используя известное из векторного анализа векторное тождество
\[
(\mathbf{u} \cdot \nabla)\mathbf{u} = \text{rot} \mathbf{u} \times \mathbf{u} + \nabla \left( \frac{u^2}{2} \right)
\]
и предполагая, что массовая сила является потенциальной, так что $\mathbf{g} = -\nabla G$, где $G$ --- потенциал силы тяжести, перепишем (5.29) в виде:
\begin{equation}\tag{5.30}
\text{rot} \mathbf{u} \times \mathbf{u} = -\nabla \left( \frac{p}{\rho} + \frac{u^2}{2} + G \right), \quad \text{div} \mathbf{u} = 0.
\end{equation}


Многие течения идеальной жидкости являются \textit{безвихревыми}, т. е. удовлетворяют условию $\text{rot} \mathbf{u} \equiv 0$. Для таких течений модель (5.30) принимает вид
\begin{equation}\tag{5.31}
\text{rot} \mathbf{u} = 0, \quad \text{div} \mathbf{u} = 0,
\end{equation}
\begin{equation}\tag{5.32}
\nabla \left( \frac{p}{\rho} + \frac{u^2}{2} + G \right) = 0 \Rightarrow \frac{p}{\rho} + \frac{u^2}{2} + G = \text{const}, \quad \rho = \text{const}.
\end{equation}


Важно отметить, что гидродинамическая модель в данном случае разделилась, т. е. свелась к системе двух уравнений (5.31) для скорости $\mathbf{u}$ и уравнению (5.32) для давления. Оно называется в гидродинамике \textit{уравнением Бернулли}.

Напомним, что условие $\text{rot} \mathbf{u} = 0$ эквивалентно (по крайней мере для односвязных областей) условию потенциальности потока, т. е. существованию такой функции $\varphi$, называемой \textit{потенциалом скорости}, что $\mathbf{u} = -\text{grad} \varphi$. Подставляя это соотношение во второе уравнение в (5.31), приходим с учетом (3.18) к \textit{уравнению Лапласа}:
\begin{equation}\tag{5.33}
\text{div}(\text{grad} \varphi) \equiv \Delta \varphi = 0
\end{equation}
для потенциала $\varphi$. Следовательно, уравнение (5.33) образует математическую модель стационарного безвихревого движения идеальной однородной несжимаемой жидкости. Таким образом, задача изучения потенциального движения идеальной несжимаемой жидкости сводится к нахождению решения уравнения Лапласа.
  \end{enumerate}

\item \textbf{Волновое уравнение для звука}
  \begin{enumerate}[label=(\alph*)]
  \item Опишите процедуру линеаризации уравнений гидродинамики. Какие допущения при этом делаются?\subsection*{2.2. Линеаризация для малых возмущений}

Предположим, что жидкость покоится ($\mathbf{u} = 0$) в состоянии равновесия с постоянными параметрами $\rho = \rho_0, p = p_0$. Рассмотрим малые возмущения этого состояния:
\[
\mathbf{u} = \mathbf{u}', \quad p = p_0 + p', \quad \rho = \rho_0 + \rho'.
\]
Подставим эти выражения в уравнения, пренебрегая величинами второго и более высокого порядка малости. Из уравнения неразрывности получаем:
\begin{equation*}
\frac{\partial \rho'}{\partial t} + \rho_0 \text{div} \mathbf{u}' = 0.
\end{equation*}
Из уравнения Эйлера (при $\mathbf{f} = 0$):
\begin{equation*}
\rho_0 \frac{\partial \mathbf{u}'}{\partial t} = -\nabla p'.
\end{equation*}

\subsection*{2.3. Уравнение состояния и скорость звука}

Разложим функцию $P(\rho)$ в ряд Тейлора в окрестности $\rho_0$:
\[
p = P(\rho_0) + \left. \frac{dP}{d\rho} \right|_{\rho=\rho_0} (\rho - \rho_0) + O((\rho - \rho_0)^2) = p_0 + c^2 \rho' + \dots,
\]
где введено обозначение скорости звука $c$:
\begin{equation}\tag{3}
c = \sqrt{\frac{dp}{d\rho}}.
\end{equation}
Для идеального газа при адиабатическом процессе ($\rho = \text{const}$) формула преобразуется к виду:
\begin{equation}\tag{5}
c = \sqrt{\gamma \frac{P}{\rho}} = \sqrt{\frac{\gamma RT}{M}},
\end{equation}
где $\gamma$ --- показатель адиабаты, $R$ --- универсальная газовая постоянная, $T$ --- абсолютная температура, $M$ --- молярная масса газа.

\subsection*{2.4. Волновое уравнение для звукового давления}

Исключим из линеаризованных уравнений $\rho'$ и $\mathbf{u}'$. Используя связь $p' = c^2 \rho'$, приходим к классическому \textit{волновому уравнению} для давления:
\begin{equation*}
\frac{\partial^2 p}{\partial t^2} = c^2 \Delta p, \quad p \equiv p'.
\end{equation*}
В общем случае при наличии внешних сил $F$ уравнение становится неоднородным:
\begin{equation*}
\frac{\partial^2 p}{\partial t^2} = c^2 \Delta p + F.
\end{equation*}

\subsection*{3. Постановка задач для волнового уравнения}

Для выделения единственного решения необходимо задать начальные и граничные условия.

\subsubsection*{3.1. Начальные условия. Задача Коши}
В неограниченном пространстве $\mathbb{R}^3$ задаются условия при $t=0$:
\begin{equation}\tag{6.16}
p|_{t=0} = p_0(\mathbf{x}), \quad \left. \frac{\partial p}{\partial t} \right|_{t=0} = p_1(\mathbf{x}), \quad \mathbf{x} \in \mathbb{R}^3.
\end{equation}

\subsubsection*{3.2. Граничные условия}
В ограниченной области $\Omega$ с границей $\Gamma$ задаются условия трех основных типов:
\begin{itemize}
    \item \textbf{Условие Дирихле} (1-го рода): задается значение давления $p|_\Gamma = g(\mathbf{x}, t)$.
    \item \textbf{Условие Неймана} (2-го рода): задается нормальная производная (поток) $\left. \frac{\partial p}{\partial n} \right|_\Gamma = g(\mathbf{x}, t)$.
    \item \textbf{Условие третьего рода} (смешанное): $ap + b \frac{\partial p}{\partial n} = g$ на $\Gamma$.
\end{itemize}

\subsubsection*{3.3. Задача излучения и дифракции}
В неограниченных областях дополнительно требуются \textit{условия излучения Зоммерфельда} на бесконечности, гарантирующие, что волна уходит на бесконечность, а не приходит из нее:
\begin{equation*}
p(\mathbf{x}) = O(|\mathbf{x}|^{-1}), \quad \frac{\partial p}{\partial |\mathbf{x}|} - ikp = o(|\mathbf{x}|^{-1}), \quad |\mathbf{x}| \to \infty.
\end{equation*}
  \item Как вводится скорость звука? Запишите волновое уравнение для звукового давления.
  \item Сформулируйте задачу Коши и начально-краевые задачи для волнового уравнения.
  \end{enumerate}

\item \textbf{Уравнения Максвелла}
  \begin{enumerate}[label=(\alph*)]
  \item Запишите систему уравнений Максвелла в дифференциальной форме. Объясните физический смысл каждого уравнения.
  \item Выведите волновое уравнение для \(E\) в непроводящей среде. Чему равна скорость электромагнитных волн?
  \item Почему Максвелл предсказал существование электромагнитных волн?
  \section*{4. Система уравнений Максвелла}

Электромагнитные процессы в среде описываются \textbf{уравнениями Максвелла}. В системе СИ они имеют вид:
\begin{align}
\text{div} \mathbf{D} &= \rho_e, \tag{7.1} \\
\text{div} \mathbf{B} &= 0, \tag{7.2} \\
\text{rot} \mathbf{E} &= -\frac{\partial \mathbf{B}}{\partial t}, \tag{7.3} \\
\text{rot} \mathbf{H} &= \frac{\partial \mathbf{D}}{\partial t} + \mathbf{J}. \tag{7.4}
\end{align}

Здесь:
\begin{itemize}
    \item $\mathbf{E}$ --- напряженность электрического поля;
    \item $\mathbf{D}$ --- электрическая индукция;
    \item $\mathbf{H}$ --- напряженность магнитного поля;
    \item $\mathbf{B}$ --- магнитная индукция;
    \item $\rho_e$ --- объемная плотность электрических зарядов;
    \item $\mathbf{J}$ --- плотность электрического тока.
\end{itemize}

\subsection*{4.1. Физический смысл уравнений}

\begin{enumerate}
    \item Уравнение (7.1) является дифференциальной формой \textbf{закона Кулона}: источники электрического поля --- заряды.
    \item Уравнение (7.2) отражает отсутствие магнитных зарядов (монополей) --- линии магнитной индукции всегда замкнуты.
    \item Уравнение (7.3) представляет \textbf{закон электромагнитной индукции Фарадея}: переменное магнитное поле порождает вихревое электрическое поле.
    \item Уравнение (7.4) обобщает \textbf{закон Ампера}: магнитное поле создается токами проводимости и токами смещения ($\partial \mathbf{D} / \partial t$).
\end{enumerate}

Для замыкания системы необходимы \textbf{материальные уравнения}, связывающие $\mathbf{D, B, J}$ с $\mathbf{E}$ и $\mathbf{H}$. В простейшем случае однородной изотропной среды:
\begin{equation*}
\mathbf{D} = \varepsilon \varepsilon_0 \mathbf{E}, \quad \mathbf{B} = \mu \mu_0 \mathbf{H}, \quad \mathbf{J} = \sigma \mathbf{E} + \mathbf{J}^{st}.
\end{equation*}
Здесь $\varepsilon_0, \mu_0$ --- электрическая и магнитная постоянные, $\varepsilon, \mu$ --- относительные диэлектрическая и магнитная проницаемости, $\sigma$ --- удельная проводимость, $\mathbf{J}^{st}$ --- плотность сторонних токов.
  \end{enumerate}

\item \textbf{Классификация уравнений второго порядка}
  \begin{enumerate}[label=(\alph*)]
  \item Что такое квазилинейное и линейное уравнение? Приведите примеры.
  \item Как строится квадратичная форма из старших коэффициентов? Как по ней определить тип уравнения в точке?
  \item Дайте определения эллиптического, гиперболического и параболического типов в области. Приведите примеры.
  \item Что такое уравнение смешанного типа? Приведите пример уравнения Трикоми и объясните, где оно меняет тип.\section*{5. Вывод волнового уравнения для электромагнитного поля}

\subsection*{5.1. Упрощения для однородной непроводящей среды}
Рассмотрим однородную непроводящую среду ($\sigma = 0$), в которой нет сторонних токов ($\mathbf{J}^{st} = 0$). Тогда $\mathbf{J} = 0$. Материальные уравнения принимают вид:
\begin{equation*}
\mathbf{D} = \varepsilon\varepsilon_0 \mathbf{E}, \quad \mathbf{B} = \mu\mu_0 \mathbf{H},
\end{equation*}
где $\varepsilon, \mu$ --- постоянные.

\subsection*{5.2. Получение волновых уравнений для $\mathbf{E}$ и $\mathbf{H}$}
Применим операцию $\text{rot}$ к уравнению Максвелла для $\mathbf{E}$ (7.3):
\begin{equation*}
\text{rot rot} \mathbf{E} = -\frac{\partial}{\partial t} \text{rot} \mathbf{B}.
\end{equation*}
С учетом $\text{rot} \mathbf{B} = \mu\mu_0 \text{rot} \mathbf{H}$ и уравнения (7.4) ($\text{rot} \mathbf{H} = \partial\mathbf{D}/\partial t = \varepsilon\varepsilon_0 \partial\mathbf{E}/\partial t$) получаем:
\begin{equation*}
\text{rot rot} \mathbf{E} = -\mu\mu_0 \frac{\partial}{\partial t} \left( \varepsilon\varepsilon_0 \frac{\partial \mathbf{E}}{\partial t} \right) = -\varepsilon\varepsilon_0\mu\mu_0 \frac{\partial^2 \mathbf{E}}{\partial t^2}.
\end{equation*}
Используя векторное тождество $\text{rot rot} \mathbf{E} = \text{grad div} \mathbf{E} - \Delta \mathbf{E}$ и учитывая отсутствие зарядов ($\text{div} \mathbf{D} = 0 \Rightarrow \text{div} \mathbf{E} = 0$), имеем:
\begin{equation*}
\text{rot rot} \mathbf{E} = -\Delta \mathbf{E}.
\end{equation*}
Подставляя, приходим к волновому уравнению для электрического поля:
\begin{equation*}
\boxed{\frac{\partial^2 \mathbf{E}}{\partial t^2} = \frac{1}{\varepsilon\varepsilon_0\mu\mu_0} \Delta \mathbf{E}}.
\end{equation*}
Аналогично получается такое же уравнение для магнитного поля $\mathbf{H}$.

\subsection*{5.3. Скорость электромагнитных волн}
Обозначим:
\begin{equation*}
a^2 = \frac{1}{\varepsilon\varepsilon_0\mu\mu_0}.
\end{equation*}
Для вакуума ($\varepsilon = \mu = 1$):
\begin{equation*}
a_0 = \frac{1}{\sqrt{\varepsilon_0\mu_0}} \approx 3 \cdot 10^8 \text{ м/с},
\end{equation*}
что совпадает со скоростью света. Это позволило Максвеллу сделать вывод о том, что свет --- это электромагнитные волны.

\section*{6. Сравнительный анализ основных типов уравнений}

Выведенные уравнения представляют три основных типа уравнений в частных производных второго порядка:

\begin{enumerate}
    \item \textbf{Эллиптический тип.} Пример: уравнение Пуассона $\Delta u = -f$. Характеризует стационарные поля, обладает свойством бесконечной гладкости внутри области и принципом максимума.
    \item \textbf{Параболический тип.} Пример: уравнение теплопроводности $\partial u / \partial t = a^2 \Delta u + f$. Описывает нестационарные процессы диффузии. Характеризуется бесконечной скоростью распространения возмущений.
    \item \textbf{Гиперболический тип.} Пример: волновое уравнение $\partial^2 u / \partial t^2 = c^2 \Delta u + f$. Описывает волновые процессы, характеризуется конечной скоростью распространения и наличием характеристических поверхностей.
\end{enumerate}

Каждый тип требует своей постановки условий: для эллиптических --- только граничные, для параболических --- одно начальное и граничные, для гиперболических --- два начальных и граничные.
  \end{enumerate}

\item \textbf{Корректность постановки задач}
  \begin{enumerate}[label=(\alph*)]
  \item Сформулируйте определение корректно поставленной задачи по Адамару. Что такое существование, единственность, устойчивость?
  \item Приведите пример корректной задачи (задача Коши для ОДУ).
  \item Опишите пример Адамара некорректной задачи (задача Коши для уравнения Лапласа). Почему она некорректна?
  \item Что такое классическое и обобщенное решение?
  \item В чем отличие корректности по Тихонову (условно-корректные задачи)?
  \end{enumerate}
\end{enumerate}
\subsection*{7. Дополнительные сведения о постановках задач}

Для выделения единственного решения уравнений второго порядка необходимо задать соответствующие дополнительные условия.

\begin{itemize}
    \item \textbf{Для эллиптических уравнений} (например, уравнение Пуассона $\Delta u = -f$): задача сводится исключительно к заданию граничных условий на границе $\Gamma$ области $\Omega$.
    \item \textbf{Для параболических уравнений} (например, уравнение теплопроводности): требуется задать одно начальное условие (состояние среды в момент $t=0$) и граничные условия на $\Gamma$ для всех $t > 0$.
    \item \textbf{Для гиперболических уравнений} (например, волновое уравнение): необходимо задать два начальных условия (начальное положение и начальную скорость, так как уравнение второго порядка по времени) и, при необходимости, граничные условия.
\end{itemize}

\subsection*{8. Сводная классификация}

Основные типы линейных дифференциальных уравнений в частных производных второго порядка можно представить в виде таблицы:

\begin{table}[h]
\centering
\begin{tabular}{|l|l|l|}
\hline
\textbf{Тип} & \textbf{Пример} & \textbf{Физический смысл} \\ \hline
Эллиптический & $\Delta u = -f$ & Стационарные процессы \\ \hline
Параболический & $u_t = a^2 \Delta u + f$ & Диффузия, теплопроводность \\ \hline
Гиперболический & $u_{tt} = c^2 \Delta u + f$ & Колебания, волны \\ \hline
\end{tabular}
\end{table}



Важно отметить, что гиперболический тип уравнений характеризуется наличием конечной скорости распространения возмущений и существованием характеристических поверхностей, в то время как для параболических уравнений скорость распространения возмущений бесконечна, что приводит к немедленному сглаживанию начальных данных.
\section*{7. Классификация уравнений в частных производных второго порядка}

Рассмотрим линейное уравнение второго порядка от двух независимых переменных $x_1, x_2$:
\begin{equation}\tag{7.1}
A_{11} \frac{\partial^2 u}{\partial x_1^2} + 2A_{12} \frac{\partial^2 u}{\partial x_1 \partial x_2} + A_{22} \frac{\partial^2 u}{\partial x_2^2} + F\left(x_1, x_2, u, \frac{\partial u}{\partial x_1}, \frac{\partial u}{\partial x_2}\right) = 0.
\end{equation}

Классификация таких уравнений основывается на знаке дискриминанта квадратичной формы, соответствующей главной части оператора (старшим производным). Определим дискриминант $D$ как:
\begin{equation}\tag{7.2}
D = A_{12}^2 - A_{11}A_{22}.
\end{equation}пз

В зависимости от значения $D$ уравнение (7.1) в данной точке $(x_1, x_2)$ классифицируется следующим образом:

\begin{itemize}
    \item \textbf{Гиперболический тип}: если $D > 0$.
    \item \textbf{Параболический тип}: если $D = 0$.
    \item \textbf{Эллиптический тип}: если $D < 0$.
\end{itemize}

\subsection*{7.1. Физическая интерпретация типов уравнений}

Для уравнений, содержащих производную по времени $t$ (где $x_1 = t$, $x_2 = x$), эта классификация тесно связана с физическими процессами, описываемыми уравнениями:

\begin{enumerate}
    \item \textbf{Гиперболические уравнения} ($D > 0$): описывают волновые процессы, где возмущение распространяется с конечной скоростью. Примером служит волновое уравнение.
    \item \textbf{Параболические уравнения} ($D = 0$): описывают процессы диффузии и теплопроводности. Здесь скорость распространения возмущения формально бесконечна, а начальные данные со временем сглаживаются.
    \item \textbf{Эллиптические уравнения} ($D < 0$): описывают стационарные процессы, находящиеся в равновесии. К ним относятся уравнение Лапласа и уравнение Пуассона.
\end{enumerate}

Данн

\section*{Часть 2. Задачи и упражнения}

\begin{enumerate}
\item \textbf{Операторы}
  \begin{enumerate}[label=(\alph*)]
  \item Дано скалярное поле \(u = x^2y + z^3\). Найдите \(\operatorname{grad} u\) в точке \((1,2,1)\).
  \item Для векторного поля \(\mathbf{v} = (xy, yz, zx)\) вычислите \(\operatorname{div}\mathbf{v}\) и \(\operatorname{rot}\mathbf{v}\).
  \end{enumerate}

\item \textbf{Уравнение Лапласа}
  \begin{enumerate}[label=(\alph*)]
  \item Проверьте, что функция \(u=\ln \frac{1}{r}\) (где \(r=\sqrt{x^2+y^2}\)) удовлетворяет уравнению Лапласа в \(\mathbb{R}^2\setminus\{0\}\).
  \item Покажите, что потенциал точечного заряда \(\varphi=\dfrac{q}{4\pi\varepsilon_0 r}\) гармоничен всюду, кроме начала координат.
  \end{enumerate}

\item \textbf{Уравнение теплопроводности}
  \begin{enumerate}[label=(\alph*)]
  \item Запишите одномерное уравнение теплопроводности для стержня с источниками. Каков физический смысл каждого члена?
  \item Сформулируйте начально-краевую задачу для стержня длины \(l\) с теплоизолированными концами (условие Неймана).
  \end{enumerate}

\item \textbf{Волновое уравнение}
  \begin{enumerate}[label=(\alph*)]
  \item Покажите, что функция \(u(x,t)=f(x-ct)+g(x+ct)\) удовлетворяет одномерному волновому уравнению.
  \item Каков физический смысл прямой и обратной волн?
  \end{enumerate}

\item \textbf{Классификация}
  \begin{enumerate}[label=(\alph*)]
  \item Определите тип следующих уравнений:
  \[
  \text{а) }u_{xx}-2u_{xy}+u_{yy}=0,\quad
  \text{б) }u_{xx}+4u_{yy}=0,
  \]
  \[
  \text{в) }u_{tt}-4u_{xx}=0,\quad
  \text{г) }u_{xx}+2u_{xy}+u_{yy}+u_x=0.
  \]
  \item Для каждого найдите матрицу коэффициентов и, если возможно, дискриминант.
  \end{enumerate}

\item \textbf{Корректность}
  \begin{enumerate}[label=(\alph*)]
  \item Приведите пример, когда решение существует, но не единственно.
  \item Почему задача Коши для уравнения Лапласа неустойчива?
  \item Как можно сделать некорректную задачу условно-корректной?
  \end{enumerate}
\end{enumerate}

\end{document}
